\chapter{Experimentos y Resultados}\label{cap4:resultados}

Este capítulo presenta los resultados de dos experimentos completos y el
diagnóstico preliminar de un tercero, ejecutados en GPU NVIDIA RTX~5050
(CUDA 12.8, PyTorch 2.11, SEED = 42). Los experimentos siguen una complejidad
creciente: validación analítica 1D (NB01), validación 2D con campo complejo
(NB02) y simulación de \textit{speckle} óptico aún en progreso (NB03).

% ──────────────────────────────────────────────────────────────────────────────
\section{NB01: Validación 1D, Helmholtz con solución analítica}
\label{sec:resultados1d}

\subsection{Configuración experimental}

El primer experimento valida la arquitectura SIREN frente a la solución analítica
exacta de la ecuación de Helmholtz 1D. Se empleó la arquitectura SIREN
$5 \times 64$ ($16{,}833$ parámetros) con los hiperparámetros de la
Tabla~\ref{tab:hiperparametros}. Los puntos de colocación son $N_c = 2{,}000$
puntos uniformes en $[0,1]$. Las condiciones de frontera Dirichlet físicas se
imponen en los extremos $x=0$ y $x=1$; adicionalmente, se incorporan tres
puntos interiores de supervisión de la solución analítica
($x=0.25,\ 0.5,\ 0.75$). Estos cinco puntos conocidos evitan el colapso hacia
la solución constante que produce la condición simétrica $E(0)=E(1)=1$.

\subsection{Resultados de precisión}

La Tabla~\ref{tab:resultados1d} presenta las métricas de precisión obtenidas
tras el entrenamiento completo Adam (15{,}000 épocas) seguido de L-BFGS
(202 evaluaciones de la función objetivo), con un tiempo total de
131.1 segundos.

\begin{table}[htbp]
\centering
\begin{threeparttable}
\caption{Métricas de precisión: PINN-SIREN 1D ($k = 2\pi$).}
\label{tab:resultados1d}
\begin{tabular}{lc}
\toprule
Métrica & Valor \\
\midrule
Error $L^2$ relativo (\%)    & \textbf{0.006} \\
MSE                          & $1.5 \times 10^{-9}$ \\
RMSE                         & $3.9 \times 10^{-5}$ \\
MAE                          & $3.50 \times 10^{-5}$ \\
Error máximo                 & $5.77 \times 10^{-5}$ \\
$R^2$                        & 1.000000 \\
Correlación de Pearson       & 1.000000 \\
\midrule
Residuo de Helmholtz, RMS    & $3.63 \times 10^{-3}$ \\
Residuo de Helmholtz, máximo & $1.48 \times 10^{-2}$ \\
Residuo relativo             & $1.30 \times 10^{-4}$ \\
\midrule
Épocas Adam                  & 15{,}000 / 15{,}000 \\
Evaluaciones L-BFGS          & 202 / 500 \\
Tiempo total (s)             & 131.1 \\
\bottomrule
\end{tabular}
\begin{tablenotes}\small
    \item Evaluado en malla uniforme de 1{,}000 puntos en $[0,1]$.
    Arquitectura SIREN: 5 capas $\times$ 64 neuronas, $\omega_0 = 1.0$.
    Solución analítica: $E_\text{exacta}(x) = \cos(kx)$.
    GPU: NVIDIA RTX 5050, SEED = 42, laptop conectada a corriente y sin
    procesos concurrentes. ``Evaluaciones L-BFGS'' cuenta las llamadas a la
    función de pérdida durante la búsqueda de línea \textit{strong Wolfe},
    no iteraciones del optimizador. En una corrida previa sin estas
    condiciones el tiempo total fue 504.2 s: el mismo hardware, ejecutando
    en régimen térmico degradado (frecuencia de reloj de GPU reducida por
    carga sostenida), reprodujo idénticamente $L^2$, $R^2$ y el resto de
    métricas de precisión, pero a 8.75 ms/época en el primer tramo y
    60.4 ms/época tras la época $\sim$8{,}700. El tiempo de entrenamiento en
    hardware portátil depende por tanto de las condiciones térmicas de la
    sesión y no debe interpretarse como una cifra de rendimiento del
    algoritmo. Una tercera ejecución, al reejecutar el notebook completo
    de forma concurrente con otro entrenamiento en la misma GPU, arrojó
    357.5 s y volvió a reproducir de forma idéntica $L^2$, MSE, RMSE, MAE,
    error máximo, $R^2$, correlación, número de épocas, evaluaciones L-BFGS
    y mejor pérdida de Adam ($4.847\times10^{-4}$). Las métricas de precisión
    son por tanto estables frente a las condiciones de ejecución; sólo el
    tiempo varía.
\end{tablenotes}
\end{threeparttable}
\end{table}

\subsection{Análisis y comparativa con el estado del arte}

El error $L^2 = 0.006\%$ se mantiene por debajo del umbral de la tesis
($< 5\%$) por casi tres órdenes de magnitud. Este resultado es atribuible a
la ventaja intrínseca de
la activación sinusoidal SIREN para representar funciones oscilatorias: la
red aprende la solución $\cos(kx)$ de forma casi exacta, con
$R^2 = 1.000000$ y correlación de Pearson $= 1.000000$ en toda la malla de
evaluación.

La Figura~\ref{fig:resultados1d} muestra la comparación entre la predicción de
la red y la solución analítica, así como la distribución del error puntual.

\begin{figure}[htbp]
    \centering
    \includegraphics[width=0.9\textwidth]{resultados_pinn_1d.png}
    \caption{Validación PINN-SIREN para Helmholtz 1D.
    Panel A: comparación entre la predicción de la red y la solución
    analítica $E_\text{exacta}(x) = \cos(kx)$ en 1{,}000 puntos uniformes en
    $[0,1]$. Panel B: error puntual $|\hat{E}(x) - E_\text{exacta}(x)|$.
    Panel C: curva de pérdida total durante las fases Adam y L-BFGS; la
    transición a la época 15{,}000 es visible como una caída abrupta.
    Error $L^2 = 0.006\%$, $R^2 = 1.000000$. GPU: NVIDIA RTX 5050, SEED = 42.}
    \label{fig:resultados1d}
\end{figure}

La Figura~\ref{fig:metricas1d} complementa este resultado con la distribución
del error y las métricas de ajuste adicionales.

\begin{figure}[htbp]
    \centering
    \includegraphics[width=\textwidth]{metricas_adicionales_1d.png}
    \caption{Métricas adicionales de entrenamiento NB01. Panel A:
    distribución del error absoluto $|\hat{E}(x) - E_\text{exacta}(x)|$.
    Panel B: error relativo puntual (\%). Panel C: dispersión de la
    predicción PINN contra la solución exacta, con la línea $y = x$ como
    referencia de ajuste perfecto.}
    \label{fig:metricas1d}
\end{figure}

\subsection{Barrido de $\omega_0$: evidencia cuantitativa del sesgo espectral}
\label{sec:omega0_sweep}

La regla de calibración $\omega_0 \approx k/(2\pi)$ (Tabla~\ref{tab:hiperparametros})
se documentó originalmente a partir de un único caso puntual: $\omega_0 = 30$
(el valor recomendado por \textcite{sitzmann2020siren} para señales de imagen)
desestabiliza el entrenamiento con $k \approx 2\pi$. Para verificar esta regla
más allá de un solo punto, se entrenó el mismo problema de NB01 con
$\omega_0 \in \{1.0,\ 5.0,\ 15.0,\ 30.0\}$ durante 3{,}000 épocas Adam
idénticas (sin L-BFGS), manteniendo $k = 2\pi$ fijo.

\begin{table}[htbp]
\centering
\begin{threeparttable}
\caption{Barrido de $\omega_0$ sobre el problema de NB01 (3{,}000 épocas Adam).}
\label{tab:omega0_sweep}
\begin{tabular}{lcc}
\toprule
$\omega_0$ & Pérdida total final & Observación \\
\midrule
1.0  & \textbf{0.2656} & Calibrado ($\omega_0 = k/2\pi$) \\
5.0  & 0.5355 & Descalibrado \\
15.0 & 0.4375 & Descalibrado \\
30.0 & 0.6030 & Descalibrado (valor de \textcite{sitzmann2020siren} para imágenes) \\
\bottomrule
\end{tabular}
\begin{tablenotes}\small
    \item Corridas nuevas y cortas, independientes del entrenamiento final de
    NB01 (15{,}000 épocas + L-BFGS, Tabla~\ref{tab:resultados1d}); no
    modifican ni reemplazan esos resultados. Mismo $N_c = 2{,}000$,
    $\lambda_\text{fís} = 1.0$ y puntos de frontera que NB01.
    Ninguna corrida diverge (pérdida infinita o NaN) dentro de las 3{,}000
    épocas: la desestabilización numérica documentada para $\omega_0 = 30$ en
    el entrenamiento completo de 15{,}000 épocas no se reproduce como
    divergencia literal en este barrido corto, pero sí se observa la
    tendencia esperada de peor convergencia a mayor descalibración.
\end{tablenotes}
\end{threeparttable}
\end{table}

El valor calibrado $\omega_0 = 1.0$ converge a una pérdida menor que
$\omega_0 = 15$ y $\omega_0 = 30$ en la misma cantidad de épocas. La
comparación con $\omega_0 = 5$ no es concluyente: una repetición independiente
del mismo barrido, ejecutada en CPU en lugar de GPU
(\texttt{results/diagnostics/nb01\_omega0\_sweep.json}), invierte el orden entre
$1$ y $5$. Lo que ambas corridas sí establecen de forma consistente es la
degradación para $\omega_0 \geq 15$; el barrido de la formulación modal
(Tabla~\ref{tab:nb03_omega0}) muestra el mismo patrón, con $1$ y $5$ próximos
entre sí y separados por más de un orden de magnitud de los valores altos. Este resultado es consistente con el sesgo espectral
descrito por \textcite{rahaman2019spectral}: la activación sinusoidal de
SIREN impone un espectro de frecuencias nativo determinado por $\omega_0$, y
cuando este no coincide con la frecuencia $k$ de la solución objetivo, la red
debe aprender a compensar ese desajuste espectral antes de ajustar la
solución misma, retrasando o degradando la convergencia.

% ──────────────────────────────────────────────────────────────────────────────
\section{NB02: Validación 2D, Helmholtz con campo complejo}
\label{sec:resultados2d}

\subsection{Configuración experimental}

El segundo experimento escala el modelo al caso bidimensional con campo complejo.
Se empleó la arquitectura SIREN $5 \times 128$ ($66{,}690$ parámetros), LHS
con $N_c = 3{,}000$ puntos de colocación y $N_b = 300$ puntos de frontera por
borde. El peso de física $\lambda_\text{fís} = 0.1$ fue seleccionado mediante
el experimento de ablación descrito en la Sección~\ref{sec:ablacion}.

\subsection{Resultados de precisión}

La Tabla~\ref{tab:resultados2d} reporta las métricas detalladas para cada
componente del campo complejo, evaluadas en una malla uniforme de
$100 \times 100$ puntos.

\begin{table}[htbp]
\centering
\begin{threeparttable}
\caption{Métricas de precisión: PINN-SIREN 2D ($k = 2\pi$, $k_x = k_y = k/\sqrt{2}$).}
\label{tab:resultados2d}
\begin{tabular}{lcc}
\toprule
Métrica & $E_\text{real}$ & $E_\text{imag}$ \\
\midrule
Error $L^2$ componente (\%) & 0.214              & 0.127              \\
MSE                      & $2.20 \times 10^{-6}$ & $8.43 \times 10^{-7}$ \\
RMSE                     & $1.48 \times 10^{-3}$ & $9.18 \times 10^{-4}$ \\
MAE                      & $1.27 \times 10^{-3}$ & $7.71 \times 10^{-4}$ \\
Error máximo             & $3.05 \times 10^{-3}$ & $2.93 \times 10^{-3}$ \\
$R^2$                    & 0.999995           & 0.999998           \\
Correlación de Pearson   & 0.999998           & 0.999999           \\
\midrule
Promedio aritmético de errores $L^2$ (\%)& \multicolumn{2}{c}{\textbf{0.171}} \\
Error $L^2$ del campo complejo (\%)      & \multicolumn{2}{c}{0.174} \\
Error $L^2$ de intensidad (\%)           & \multicolumn{2}{c}{0.248} \\
RMSE de fase circular (rad)              & \multicolumn{2}{c}{$1.23 \times 10^{-3}$} \\
Residuo de Helmholtz, RMS relativo       & \multicolumn{2}{c}{$1.96 \times 10^{-4}$} \\
\bottomrule
\end{tabular}
\begin{tablenotes}\small
    \item Evaluado en malla uniforme de $100 \times 100$ puntos.
    Arquitectura: SIREN $5 \times 128$, $\omega_0 = 1.0$, 66{,}690 parámetros.
    Puntos de colocación: LHS $N_c = 3{,}000$.
    Frontera: onda plana $e^{i(k_x x + k_y y)}$ en 4 bordes, $N_b = 300$ pts/borde.
    Entrenamiento: Adam (8{,}737 épocas) + L-BFGS (1{,}035 evaluaciones de la
    función objetivo), tiempo total 209.9 s (GPU RTX 5050 laptop, conectada
    a corriente, sin procesos concurrentes).
    Dispersión entre semillas (SEED $\in \{42, 123, 777\}$):
    error $L^2$ promedio $= 0.192 \pm 0.089$\% (media $\pm$ desviación
    estándar poblacional, $\mathrm{ddof}=0$).
\end{tablenotes}
\end{threeparttable}
\end{table}

\subsection{Análisis de resultados}

El error $L^2$ promedio de $0.171\%$ se mantiene por debajo del umbral de la
tesis ($< 5\%$) por más de un orden de magnitud. La precisión se mantiene alta
en ambas componentes del campo complejo, con $R^2 > 0.9999$ para
$E_\text{real}$ y $E_\text{imag}$.

En un experimento preliminar con $d = 64$ (mismo $N_c$, $N_b$ y protocolo de
optimización que la Tabla~\ref{tab:hiperparametros}) el error promedio fue
$0.436\%$, frente a $0.171\%$ con $d = 128$. Esta comparación involucra
además una diferencia en el número de evaluaciones de L-BFGS entre ambas
corridas, por lo que la mejora de $2.5\times$ no puede atribuirse de forma
aislada al aumento de capacidad de la red; se reporta como evidencia
cualitativa de que $d = 128$ es preferible, no como una ablación controlada.
No se dispone de un experimento que aísle la contribución específica del
Muestreo Latino Hipercúbico frente a una malla cartesiana equivalente.

La Figura~\ref{fig:resultados2d} muestra los mapas de calor de la solución
exacta y la predicción PINN-SIREN para ambas componentes del campo.

\begin{figure}[htbp]
    \centering
    \includegraphics[width=0.9\textwidth]{resultados_pinn_2d.png}
    \caption{Validación PINN-SIREN para Helmholtz 2D con campo complejo.
    Fila superior (A1-A3): solución analítica, predicción PINN-SIREN y error
    absoluto para $E_\text{real}$. Fila central (B1-B3): lo mismo para
    $E_\text{imag}$. Fila inferior: intensidad predicha $I = |E|^2$ (C1),
    puntos de colocación LHS (C2) y curvas de pérdida durante Adam y L-BFGS
    (C3). La diferencia visual entre solución exacta y predicción es muy
    pequeña a esta escala; los errores puntuales son del orden de $10^{-3}$
    (véase la Tabla~\ref{tab:resultados2d}). Malla de evaluación
    $100 \times 100$, SEED = 42.}
    \label{fig:resultados2d}
\end{figure}

La Figura~\ref{fig:metricas2d} complementa este resultado con la distribución
espacial del error puntual para ambas componentes del campo.

\begin{figure}[htbp]
    \centering
    \includegraphics[width=\textwidth]{metricas_adicionales_2d.png}
    \caption{Métricas adicionales de entrenamiento NB02. Paneles A-B:
    distribución del error absoluto $|E_\text{real} - \hat{E}_\text{real}|$
    y $|E_\text{imag} - \hat{E}_\text{imag}|$. Paneles C, F: dispersión de
    la predicción PINN contra la solución exacta para cada componente.
    Paneles D-E: error relativo puntual (\%) en la malla $100 \times 100$;
    los valores más altos se concentran en bandas antidiagonales
    ($x + y \approx$ const.) donde la solución exacta cruza por cero, ya que
    el error relativo diverge al dividir por un valor cercano a cero, y no en
    las esquinas del dominio. La curva de pérdida de entrenamiento se
    muestra en el Panel C3 de la Figura~\ref{fig:resultados2d}.}
    \label{fig:metricas2d}
\end{figure}

\subsection{Reproducibilidad multi-semilla}

Para evaluar la robustez estadística del método, el experimento NB02 se repitió
con tres semillas diferentes: SEED $\in \{42, 123, 777\}$. Los resultados se
presentan en la Tabla~\ref{tab:multiseed}.

\begin{table}[htbp]
\centering
\begin{threeparttable}
\caption{Reproducibilidad multi-semilla: NB02 ($\lambda_\text{fís} = 0.1$).}
\label{tab:multiseed}
\begin{tabular}{lccc}
\toprule
SEED & Error $L^2$ prom. (\%) & Épocas Adam & Tiempo (s) \\
\midrule
42   & 0.1707 & 8{,}737  & 220.7 \\
123  & 0.0949 & 9{,}441  & 225.9 \\
777  & 0.3100 & 10{,}443 & 253.2 \\
\midrule
Media $\pm$ Desv. est. & $0.192 \pm 0.089$ & - & - \\
\bottomrule
\end{tabular}
\begin{tablenotes}\small
    \item Misma arquitectura, hiperparámetros y protocolo de entrenamiento
    (\texttt{scripts/experiments/run\_multiseed.py}), ejecutadas en la misma
    sesión para que los tiempos sean comparables entre sí. GPU: NVIDIA RTX
    5050, laptop conectada a corriente. La semilla 42 reproduce el resultado
    de NB02 ($L^2 = 0.1707\%$ vs.\ $0.171\%$, dentro del redondeo).
\end{tablenotes}
\end{threeparttable}
\end{table}

La desviación estándar inter-semilla de $\pm 0.089\%$ representa un
coeficiente de variación del 46\% respecto a la media, considerable en
términos relativos, pero las tres semillas se mantienen entre 16 y 53 veces
por debajo del umbral de la tesis ($< 5\%$): la inicialización aleatoria
afecta la precisión final de forma apreciable, sin comprometer el
cumplimiento de la hipótesis (i) en ninguno de los tres casos.

% ──────────────────────────────────────────────────────────────────────────────
\subsection{Ablación de profundidad: 4 frente a 5 capas}
\label{sec:ablacion_capas}

La profundidad de $5$ capas ocultas se heredó de NB01 sin contrastarla. Para
verificar que no es una elección arbitraria se repitieron ambos experimentos con
$4$ capas, manteniendo el resto del protocolo idéntico.

\begin{table}[htbp]
\centering
\begin{threeparttable}
\caption{Ablación de profundidad en NB01 y NB02.}
\label{tab:ablacion_capas}
\begin{tabular}{lccc}
\toprule
Experimento & Capas & Parámetros & Error $L^2$ (\%) \\
\midrule
NB01 (1D) & 4 & 12{,}673 & 0.0133 \\
NB01 (1D) & \textbf{5} & 16{,}833 & \textbf{0.006} \\
\midrule
NB02 (2D) & 4 & 50{,}178 & 0.2279 \\
NB02 (2D) & \textbf{5} & 66{,}690 & \textbf{0.171} \\
\bottomrule
\end{tabular}
\begin{tablenotes}\small
    \item Mismo $\omega_0$, semilla, muestreo, $\lambda_\text{fís}$ y
    presupuesto de épocas que las corridas de referencia. La configuración de
    $5$ capas obtiene menor error en ambos casos, con un incremento de
    parámetros del $33\%$ en 1D y del $33\%$ en 2D. La diferencia es de un
    factor $2.2$ en 1D y $1.3$ en 2D, ambas muy por debajo del umbral del
    $5\%$: la elección de profundidad no condiciona el cumplimiento de la
    hipótesis, pero queda justificada por evidencia y no por herencia. Fuente:
    \texttt{explorations/architecture\_ablation\_4layers/output/}.
\end{tablenotes}
\end{threeparttable}
\end{table}

% ──────────────────────────────────────────────────────────────────────────────
\section{Ablación del peso de física $\lambda_\text{fís}$}
\label{sec:ablacion}

Para calibrar el balance entre el término de datos y el término físico en la
función de pérdida~\eqref{eq:loss_2d}, se realizó un experimento de ablación
con $\lambda_\text{fís} \in \{0.01, 0.1, 1.0\}$, manteniendo fija la
arquitectura SIREN $5 \times 128$ y SEED = 42. Los resultados se presentan en
la Tabla~\ref{tab:ablacion_lambda}.

\begin{table}[htbp]
\centering
\begin{threeparttable}
\caption{Ablación del peso de física $\lambda_\text{fís}$: NB02.}
\label{tab:ablacion_lambda}
\begin{tabular}{lccc}
\toprule
$\lambda_\text{fís}$ & Error $L^2$ prom. (\%) & Convergencia & Observación \\
\midrule
0.01  & 0.163  & Sí & Física subponderada \\
0.1   & 0.171  & Sí & \textbf{Configuración seleccionada} \\
1.0   & \textemdash & No (\textit{crash}) & Física sobreponderada \\
\bottomrule
\end{tabular}
\begin{tablenotes}\small
    \item Arquitectura SIREN $5 \times 128$, $\omega_0 = 1.0$, SEED = 42.
    Mismo protocolo Adam + L-BFGS que NB02.
    El caso $\lambda = 1.0$ no produjo un valor de $L^2$: la corrida se
    interrumpió por una excepción de CUDA antes de completar el
    entrenamiento. El registro de la corrida asocia el fallo con el crecimiento
    del gradiente en \texttt{float32}; como no se guardó una medición independiente
    de su magnitud, se reporta como diagnóstico de la falla y no como una
    relación causal demostrada. No es un caso de no convergencia por agotar
    épocas, sino una interrupción numérica (\texttt{results/ablation\_lambda.json}).
\end{tablenotes}
\end{threeparttable}
\end{table}

El valor $\lambda_\text{fís} = 0.1$ se seleccionó por la estabilidad numérica
observada, no por precisión: el modelo con $\lambda = 0.01$ obtiene un error
ligeramente menor (0.163\% frente a 0.171\%) al asignar menor peso a la
restricción física, mientras que $\lambda = 1.0$ se interrumpió antes de
completar el entrenamiento. Entre las dos configuraciones que sí convergen,
$\lambda = 0.1$ se prefiere porque impone una restricción física más fuerte sin
un costo de precisión relevante frente al umbral de la tesis.

Esta fragilidad frente al peso relativo de los términos de pérdida no es un
fenómeno aislado de esta implementación: \textcite{krishnapriyan2021characterizing}
documentan sistemáticamente que las PINN con pesos fijos entre el término de
datos/condición inicial y el término de residuo físico pueden fallar en
converger a la solución correcta incluso cuando la función de pérdida global
parece decrecer, atribuyendo el problema a un paisaje de optimización mal
condicionado por la competencia entre ambos términos. \textcite{wang2021failure}
describen un mecanismo relacionado: los gradientes del término físico pueden
dominar en magnitud a los del término de datos durante el entrenamiento,
generando un flujo de gradiente patológico que dificulta la convergencia. La
inestabilidad observada con $\lambda = 1.0$ en este trabajo es consistente
con ambas caracterizaciones: los dos residuos (real e imaginario) duplican el
peso efectivo de la contribución del término físico, un posible factor del
mismo tipo de desbalance; esta interpretación no se aisló experimentalmente en
la corrida de ablación.

\subsection{Explicación cuantitativa: análisis del núcleo tangente neuronal (NTK)}
\label{sec:ntk_analisis}

La teoría del núcleo tangente neuronal (\textit{Neural Tangent Kernel}, NTK)
de \textcite{wang2022ntk} permite cuantificar directamente el desbalance
descrito arriba: cada término de la función de pérdida induce su propio
kernel NTK, y la magnitud relativa de sus autovalores predice qué término
domina el gradiente durante el entrenamiento. Sobre el modelo final de NB02
ya entrenado ($\lambda_\text{fís} = 0.1$), se calculó el NTK empírico del
término de datos (frontera, 80 puntos) y del término de física (residuo de
Helmholtz, 80 puntos de colocación LHS).

\begin{table}[htbp]
\centering
\begin{threeparttable}
\caption{Espectro NTK: término de datos vs. término de física (NB02).}
\label{tab:ntk}
\begin{tabular}{lcc}
\toprule
 & Traza NTK & Autovalor máximo \\
\midrule
Datos (frontera)            & $1.45 \times 10^{4}$ & $5.29 \times 10^{3}$ \\
Física (residuo Helmholtz)  & $1.24 \times 10^{7}$ & $8.49 \times 10^{6}$ \\
\midrule
Razón física / datos        & \textbf{858$\times$} & \textbf{1{,}605$\times$} \\
\bottomrule
\end{tabular}
\begin{tablenotes}\small
    \item Calculado sobre los pesos finales guardados de NB02
    (\texttt{results/models/nb02\_helmholtz2d\_gpu.pt}), sin reentrenar ni
    modificar los resultados de la Tabla~\ref{tab:resultados2d}. NTK
    aproximado con 80 puntos por término (Jacobiano de $E_\text{real}$ para
    datos, del residuo real de Helmholtz para física, respecto a los
    $66{,}690$ parámetros del modelo).
\end{tablenotes}
\end{threeparttable}
\end{table}

El espectro NTK del término de física es entre 858 y 1{,}605 veces mayor que
el del término de datos, lo que explica cuantitativamente por qué
$\lambda_\text{fís} = 1.0$ (peso igual entre ambos términos) resulta en que
el gradiente de la física domine por completo el paso de optimización,
consistente con la inestabilidad numérica reportada arriba. Esta razón no
debe interpretarse como que $\lambda_\text{fís} = 0.1$ sea el valor que
compensa exactamente el desbalance (una compensación literal sugeriría un
factor de reescalamiento considerablemente menor); el análisis NTK aquí
presentado explica la \textit{dirección y el orden de magnitud} del
desbalance identificado empíricamente en la ablación, no deriva
analíticamente el valor óptimo de $\lambda_\text{fís}$.

% ──────────────────────────────────────────────────────────────────────────────
\section{NB03: Simulación de speckle óptico}
\label{sec:resultados_speckle}

\subsection{Diagnóstico del planteamiento: por qué la formulación directa no extiende}
\label{sec:nb03_diagnostico}

La extensión directa de la arquitectura de NB02 al problema de \textit{speckle}
--una red que recibe $(\tilde x,\tilde z)$ y devuelve el campo complejo, con la
pantalla de fase impuesta como término de datos-- no converge, y la causa no es
la capacidad de la red. En la prueba de control correspondiente, el campo
aprendido reproduce la condición de frontera con un error máximo de
$9\times10^{-8}$ y alcanza un residuo de Helmholtz varios órdenes de magnitud
menor que el de cualquier corrida aceptada, y aun así difiere de la referencia
en un $92.75\%$. Las tres observaciones son simultáneamente ciertas porque el problema,
tal como estaba planteado, no tiene solución única: con datos impuestos
únicamente en la cara de entrada y sin condición de radiación, si $E$ es solución
también lo es $E$ más cualquier campo que se anule en esa frontera.

El problema pertenece por tanto a la clase de Cauchy para un operador elíptico,
mal planteada en el sentido de Hadamard. Descomponiendo en modos transversales,
la ecuación para cada coeficiente es
\begin{equation}
    a_m''(\tilde z) + (k^2-k_{x,m}^2)\,a_m(\tilde z) = 0 ,
    \label{eq:edo_modal}
\end{equation}
cuyas soluciones son oscilatorias y acotadas cuando $k_{x,m}^2<k^2$ (modos
propagantes) pero exponenciales $e^{\pm|\kappa_m|\tilde z}$ cuando
$k_{x,m}^2>k^2$ (modos evanescentes). La rama creciente amplifica cualquier error
de los datos, de modo que la dependencia continua se pierde en el subespacio
evanescente.

\subsection{Experimento puente: aislamiento del planteamiento}
\label{sec:nb02b_puente}

Para separar el efecto del planteamiento del efecto del método se construyó un
experimento con solución analítica conocida, en el que la geometría, la solución
exacta, la arquitectura ($5\times128$, $\omega_0=\sqrt{q}$) y la semilla se
mantienen fijos y sólo cambia la condición de frontera. El campo de prueba es
\begin{equation}
    E(\tilde x,\tilde y)=\sum_{n=-1}^{1} c_n\,
    e^{\,i\left(2\pi n\,\tilde x + 2\pi\sqrt{q-n^{2}}\;\tilde y\right)},
    \qquad q=1.625,
    \label{eq:campo_puente}
\end{equation}
que satisface Helmholtz de forma exacta con $k=2\pi\sqrt{q}$, es periódico en
$\tilde x$ sobre $[0,1]$ y tiene sus tres modos propagantes. El valor
$k^2=6.5\pi^2$ se eligió de modo que no coincida con ningún autovalor del
laplaciano de Dirichlet en $[0,1]^2$ --que son $\pi^2(m^2+n^2)$ con $m,n$
enteros positivos--, condición necesaria para que el caso de control tenga
solución única.

\begin{table}[htbp]
\centering
\begin{threeparttable}
\caption{Experimento puente: misma solución analítica, tres planteamientos.}
\label{tab:nb02b}
\begin{tabular}{llcc}
\toprule
Caso & Datos y formulación & $L^2$ global & $E_\text{evan}$ en $\tilde y=1$ \\
\midrule
A & Dirichlet en los 4 lados, formulación directa & \textbf{0.077\%} & $1\times10^{-5}$\% \\
B & Cauchy en $\tilde y=0$, formulación directa   & 45.14\%          & $7.39\%$ \\
C & Cauchy idéntico a B, formulación modal        & \textbf{0.155\%} & $0\%$ \\
\bottomrule
\end{tabular}
\begin{tablenotes}\small
    \item Los casos B y C reciben exactamente la misma información: valor y
    derivada normal del campo exacto en $\tilde y=0$. $E_\text{evan}$ es la
    fracción de energía de la predicción que cae fuera del cono propagante
    $|k_x|\leq k$; la solución física tiene cero por construcción. En el caso B
    esa fracción crece de $4\times10^{-4}\%$ en $\tilde y=0$ a un máximo de
    $13.2\%$, y el error local pasa de $0.209\%$ a $73.7\%$: el ajuste de
    los datos es correcto y la degradación ocurre al propagar.
    Fuente: \texttt{explorations/nb02b\_bien\_vs\_mal\_puesto/}.
\end{tablenotes}
\end{threeparttable}
\end{table}

Con información idéntica, la representación modal es unos dos órdenes de
magnitud más precisa que la directa. La diferencia no procede de los datos
disponibles ni de la arquitectura, sino del espacio de funciones representable.

\subsection{Formulación modal adoptada}
\label{sec:nb03_config}

El dominio es $\Omega_{03}=[-10,10]\times[0,1]$ en $(\tilde x,\tilde z)$, con
$1{,}024$ puntos transversales y $101$ planos longitudinales. La condición de
entrada es una pantalla de fase aleatoria periódica y correlacionada, con
longitud de correlación $\ell_\phi=0.10\lambda$ y desviación estándar $2$ rad.
De los $1{,}024$ modos transversales, $41$ son propagantes. La fracción de
energía espectral que concentran depende de la realización: $40.99\%$ en la
pantalla 42, y entre $34.55\%$ y $41.28\%$ en las cinco empleadas.

Como las fronteras laterales son periódicas, el campo admite el desarrollo
exacto $u(\tilde x,\tilde z)=\sum_m a_m(\tilde z)e^{i k_{x,m}\tilde x}$, y
Helmholtz 2D se separa en las $41$ ecuaciones diferenciales ordinarias
desacopladas de la ecuación~\eqref{eq:edo_modal}. La base se trunca al conjunto
propagante. Este truncamiento no es una simplificación de conveniencia: es la
regularización del problema de Cauchy para Helmholtz analizada por
\textcite{nguyen2016truncation}, quienes establecen la solución regularizada
acotando la frecuencia mediante un parámetro de regularización y demuestran
estabilidad y convergencia con estimaciones de error en norma $L^2$.

La condición de Cauchy se impone de forma exacta en la arquitectura,
\begin{equation}
    a_m(\tilde z)=a_m(0)+\tilde z\,a_m'(0)+\tilde z^{\,2} N_\theta(\tilde z),
    \label{eq:cauchy_dura}
\end{equation}
de modo que el valor y la derivada en $\tilde z=0$ se cumplen con independencia
de los parámetros de la red. El error es nulo por construcción algebraica; la
verificación numérica registrada da $1.28\times10^{-8}$ en RMSE, residuo de la
aritmética en \texttt{float32} y no del ajuste. Una ecuación diferencial de segundo orden con valor y
derivada fijados en un punto tiene solución única, por lo que la
ecuación~\eqref{eq:cauchy_dura} restaura la unicidad perdida en
§\ref{sec:nb03_diagnostico}. Esta estrategia de imponer la frontera dentro de la
función de prueba, en lugar de penalizarla en la pérdida, corresponde al enfoque
descrito por \textcite{sukumar2022exact}.

Dos consecuencias metodológicas se siguen de la ecuación~\eqref{eq:cauchy_dura}.
Primero, la pérdida contiene únicamente el residuo modal: no hay término de
datos y por tanto \textbf{no existe un peso $\lambda_\text{fís}$ que calibrar} en
NB03, a diferencia de NB01 y NB02. Segundo, cada salida se escala por
$s_m=k^2\max\left(|a_m(0)|,\,|a_m'(0)|/k,\,\epsilon\right)$, lo que evita que
los modos de mayor amplitud dominen el residuo y los débiles queden sin ajustar.

\begin{table}[htbp]
\centering
\begin{threeparttable}
\caption{Configuración de NB03 frente a NB01 y NB02.}
\label{tab:nb03_config}
\begin{tabular}{lll}
\toprule
Elemento & NB01 / NB02 & NB03 modal \\
\midrule
Entrada de la red      & $\tilde x$ / $(\tilde x,\tilde y)$ & $\tilde z$ escalar \\
Salidas                & $1$ / $2$ & $82$ ($41$ modos complejos) \\
Capas ocultas          & $5\times64$ / $5\times128$ & $4\times128$ \\
$\omega_0$ primera capa& $1.0$ & $1.0$ \\
Frontera               & término de datos con $\lambda_\text{fís}$ & dura, en la arquitectura \\
$\lambda_\text{fís}$   & $1.0$ / $0.1$ & no aplica \\
Optimización           & Adam $+$ L-BFGS & Adam $5{,}000$ $+$ Adam $3{,}000$ \\
Tasa de aprendizaje    & $10^{-3}$ & $2\times10^{-4}$ y $5\times10^{-5}$ \\
Colocación             & uniforme / LHS 2D & $256$ posiciones de $\tilde z$ por época \\
\bottomrule
\end{tabular}
\begin{tablenotes}\small
    \item NB03 no emplea L-BFGS en el protocolo de $\tilde z=1\lambda$ ni tiene
    peso de física que calibrar, por lo que las filas correspondientes de la
    Tabla~\ref{tab:hiperparametros} no le son aplicables. La inicialización
    sinusoidal de \textcite{sitzmann2020siren} y la autodiferenciación de
    segundo orden sí se conservan desde NB01 y NB02.
\end{tablenotes}
\end{threeparttable}
\end{table}

\subsection{Verificación de la referencia}
\label{sec:nb03_referencia}

La referencia se calcula por el método del espectro angular, que propaga cada
modo como $a_m(0)e^{i k_{z,m}\tilde z}$. El método se validó previamente contra
una propagación exacta, con un error relativo del orden de la precisión de
máquina ($\sim10^{-16}$). Las cinco
referencias empleadas se regeneraron desde cero y coinciden \textbf{bit a bit}
con las originales: la diferencia máxima es $0.000\times10^{0}$ en los ocho
campos almacenados (fase, campo objetivo, campo en todos los planos, intensidad,
intensidad del conjunto, coordenadas y espectro), para las cinco semillas.

Debe señalarse que esta referencia es \textbf{semianalítica} y no constituye un
método numérico independiente en el sentido del elemento finito: propaga cada
modo mediante la solución cerrada de la misma ecuación diferencial ordinaria que
la red minimiza. Mide, por tanto, la exactitud del solucionador frente a la
solución exacta de sus propias ecuaciones, igual que NB01 se valida contra
$\cos(2\pi\tilde x)$ y NB02 contra una onda plana. La independencia metodológica
plena corresponde a la comparación con FEM planteada como trabajo pendiente.

\subsection{Calibración de $\omega_0$ en el régimen modal}
\label{sec:nb03_omega}

La red modal recibe $\tilde z$ normalizado y cada coeficiente oscila como
$a_m(\tilde z)\propto e^{i k_{z,m}\tilde z}$ con $k_{z,m}\leq k=2\pi$: como
máximo una oscilación completa en todo el dominio. La regla
$\omega_0\approx k/2\pi$ establecida en §\ref{sec:omega0_sweep} predice por
tanto $\omega_0=1$. Para contrastarla en este segundo régimen se repitió el
barrido sobre la pantalla $42$ con protocolo idéntico.

\begin{table}[htbp]
\centering
\begin{threeparttable}
\caption{Barrido de $\omega_0$ en la formulación modal (pantalla 42, $5{,}000$ épocas Adam).}
\label{tab:nb03_omega0}
\begin{tabular}{lccc}
\toprule
$\omega_0$ & $L^2$ total & $L^2$ propagante & Residuo de la EDO \\
\midrule
1.0  & 3.819\% & \textbf{0.090\%} & $3.59\times10^{-3}$ \\
5.0  & 3.819\% & 0.081\%          & $2.61\times10^{-3}$ \\
15.0 & 3.833\% & 0.329\%          & $1.33\times10^{-2}$ \\
30.0 & 4.962\% & 3.172\%          & $5.49\times10^{-2}$ \\
\bottomrule
\end{tabular}
\begin{tablenotes}\small
    \item Mismo protocolo en los cuatro casos, sin refinamiento posterior. El
    valor $\omega_0=30$ es el recomendado por \textcite{sitzmann2020siren} para
    señales de imagen. El error propagante --la componente que la
    representación puede expresar-- se degrada en un factor de $35$ entre el
    valor calibrado y el descalibrado, y el residuo de la ecuación en un factor
    de $15$. El comportamiento es el esperado del sesgo espectral descrito por
    \textcite{rahaman2019spectral} y replica, en un problema con entrada,
    salida y física distintas, el resultado de la Tabla~\ref{tab:omega0_sweep}.
    Fuente: \texttt{results/nb03\_modal\_omega0\_sweep\_summary.json}, generado
    por \texttt{scripts/experiments/nb03\_omega0\_sweep\_summary.py}, que
    recalcula la descomposición del error desde los campos guardados.
\end{tablenotes}
\end{threeparttable}
\end{table}

Consolidado sobre las cinco pantallas, el cambio de $\omega_0=30$ a
$\omega_0=1$ reduce el error propagante medio a lo largo de la propagación de
$2.983\%$ a $0.042\%$, y el máximo de $4.791\%$ a $0.084\%$. Este
resultado es la evidencia más directa de que NB01 no es un preámbulo
pedagógico: la regla derivada allí es la que determina el desempeño de la
formulación modal.

\subsection{Resultados en \texorpdfstring{$\tilde z=1\lambda$}{z=1 lambda}}
\label{sec:nb03_resultados_z1}

\begin{table}[htbp]
\centering
\begin{threeparttable}
\caption{NB03 modal con $\omega_0$ calibrado: cinco pantallas de fase independientes en $\tilde z=1\lambda$.}
\label{tab:nb03_cinco}
\begin{tabular}{lcccccc}
\toprule
Pantalla & $L^2$ complejo & Piso evan. & Coherencia & $C_\text{ref}$ & $C_\text{PINN}$ & $|\Delta C|$ \\
\midrule
42   & 3.819\% & 3.818\% & 0.999271 & 0.9063 & 0.9012 & 0.0051 \\
123  & 2.402\% & 2.401\% & 0.999712 & 0.8827 & 0.8821 & 0.0006 \\
321  & 3.573\% & 3.573\% & 0.999362 & 0.9200 & 0.9231 & 0.0030 \\
777  & 2.947\% & 2.947\% & 0.999566 & 1.2246 & 1.2403 & 0.0157 \\
2026 & 3.064\% & 3.063\% & 0.999531 & 0.7812 & 0.7841 & 0.0029 \\
\midrule
Media & \textbf{3.161\%} & 3.160\% & 0.999488 & --- & --- & 0.0055 \\
\bottomrule
\end{tabular}
\begin{tablenotes}\small
    \item Semilla de red fija en $42$; las cinco pantallas son realizaciones de
    fase independientes. Las cinco cumplen el criterio $L^2<5\%$. El residuo
    modal RMSE medio es $2.34\times10^{-3}$ y el error de las condiciones de
    Cauchy es de $1.28\times10^{-8}$ en RMSE, nulo salvo por la aritmética en
    \texttt{float32}. Los cinco modelos se entrenaron desde cero en el
    equipo de este trabajo y reproducen corridas previas dentro de $0.08$ puntos
    porcentuales. Fuente:
    \texttt{results/nb03\_modal\_multiseed\_z1\_omega1\_summary.json}.
\end{tablenotes}
\end{threeparttable}
\end{table}

La columna del piso evanescente permite interpretar correctamente la métrica, y la Figura~\ref{fig:nb03_modal} la resume.
Pantalla por pantalla, el error total del modelo coincide con el piso hasta la
tercera cifra decimal ($3.819$ frente a $3.818$; $2.402$ frente a $2.401$;
$3.573$ frente a $3.573$), y la desviación de ambas columnas es la misma
($\pm0.497$ puntos porcentuales). Es decir: el $L^2$ total está dominado por la
energía evanescente de la referencia, que la base modal no puede representar por
construcción, y no por el ajuste de la red. Medido sobre el subespacio
representable, el error propagante medio a lo largo de la propagación es de
$0.042\%$ y el máximo de $0.084\%$.

\subsection{Dependencia del piso evanescente con la distancia}
\label{sec:nb03_piso}

La fracción evanescente decae como $e^{-|\kappa_m|\tilde z}$, de modo que el
piso es una función de la distancia de propagación y no una constante del
método.

\begin{table}[htbp]
\centering
\begin{threeparttable}
\caption{Piso evanescente y contraste de la referencia en función de la distancia.}
\label{tab:nb03_piso}
\begin{tabular}{lcc}
\toprule
$\tilde z/\lambda$ & Piso evanescente & $C$ de la referencia \\
\midrule
0.25 & 25.49\% & 0.7968 \\
0.50 & 12.61\% & 0.8283 \\
1.00 & 3.818\% & 0.9063 \\
2.00 & 0.456\% & 1.1595 \\
3.00 & 0.0598\% & 1.2069 \\
5.00 & 0.0011\% & 1.0476 \\
\bottomrule
\end{tabular}
\begin{tablenotes}\small
    \item Calculado analíticamente desde el espectro de la pantalla 42, sin
    entrenar ningún modelo. La distancia $\tilde z=1\lambda$ corresponde a
    campo cercano estricto: es la peor distancia posible para la métrica de
    campo completo, porque casi el $4\%$ de la energía es irrepresentable por
    la base propagante y el patrón aún no ha madurado estadísticamente. El
    contraste de la referencia crece monótonamente desde $C=0$ exactamente en
    $\tilde z=0$ --donde la pantalla es de fase pura y la intensidad es
    uniforme-- y alcanza el régimen de Goodman a partir de $\tilde z\geq2\lambda$.
\end{tablenotes}
\end{threeparttable}
\end{table}

\begin{figure}[htbp]
    \centering
    \includegraphics[width=\textwidth]{nb03_resultados_modal.png}
    \caption[Resultados de NB03: calibración de $\omega_0$ y piso evanescente]{
    Resultados de NB03 modal. \textbf{(a)} Error $L^2$ propagante a lo largo de
    la propagación para las cinco pantallas, con $\omega_0=30$ (rojo) y
    $\omega_0=1$ (azul); las curvas tenues son las pantallas individuales y las
    gruesas su promedio. La separación es de aproximadamente dos órdenes de
    magnitud en todo el intervalo. \textbf{(b)} Fracción de energía
    irrepresentable por la base modal --el piso evanescente-- frente a la
    distancia de propagación, en escala logarítmica, junto al error medio del
    modelo calibrado. Cómo leerla: en $\tilde z=1\lambda$ el piso es
    $92$ veces mayor que el error del modelo, de modo que la métrica de campo
    completo mide la física irrepresentable y no el ajuste de la red; ambas
    curvas se cruzan en $\tilde z\approx3.2\lambda$, distancia a partir de la
    cual el error total sí queda dominado por el modelo. Fuentes:
    \texttt{results/nb03\_modal\_multiseed\_z1\_omega1\_l2\_curves.npz} y
    \texttt{results/nb03\_angular\_spectrum\_reference\_z1\_corr0.10.npz}.}
    \label{fig:nb03_modal}
\end{figure}

Este comportamiento explica el valor $C=0.906$ de la pantalla 42 en
$\tilde z=1\lambda$ sin recurrir a una deficiencia del modelo: el patrón está
todavía en formación. El contraste transversal medido sobre la referencia crece
de $0$ en la pantalla a $0.797$ en $\tilde z=0.25\lambda$, $0.828$ en
$0.5\lambda$ y $0.906$ en $1\lambda$, mientras el tamaño de grano transversal
pasa de $0.298\lambda$ a $0.378\lambda$.

\subsection{Criterio estadístico: conjunto frente a realización individual}
\label{sec:nb03_goodman}

El criterio $|C-1|<0.1$ de \textcite{goodman2007speckle} caracteriza el
\textit{speckle} completamente desarrollado como propiedad del conjunto de
realizaciones. Aplicado a una pantalla finita individual, el criterio no se
cumple ni siquiera para la referencia exacta: su contraste varía entre $0.7812$
y $1.2246$ según la realización (Tabla~\ref{tab:nb03_cinco}). La evaluación se
realiza por ello en dos niveles:

\begin{itemize}
    \item \textbf{Conjunto.} Reuniendo las cinco pantallas, el contraste del
    campo de la PINN es $C=0.9489$, esto es $|C-1|=0.051$, dentro del criterio.
    Las referencias por espectro angular, cada una sobre $64$ realizaciones en
    $\tilde z=1\lambda$, dan contrastes de $0.9759$, $0.9757$, $0.9531$,
    $0.9619$ y $0.9637$, con media $0.9661$: el conjunto sí se encuentra en el
    régimen de \textit{speckle} completamente desarrollado.
    \item \textbf{Realización individual.} Fidelidad al contraste de su propia
    referencia, $|C_\text{PINN}-C_\text{ref}|<0.05$, que se cumple en las cinco
    pantallas con un error máximo de $0.0157$.
\end{itemize}

Esta distinción no relaja el criterio: lo aplica al objeto estadístico para el
que fue definido.

Debe reportarse una limitación asociada, que afecta a la referencia exacta y
no sólo al modelo. El test de Kolmogorov--Smirnov de bondad de ajuste a la
distribución exponencial, aplicado a los $1{,}024$ valores de intensidad de una
realización individual \emph{de la referencia por espectro angular}, sólo
alcanza $p>0.05$ en una de las cinco pantallas (valores de $p$: $0.1423$, $0.0266$, $0.0048$, $0.0000$ y
$0.0000$). El motivo es que la prueba supone muestras independientes, mientras
que los valores de intensidad de un corte espacial de \textit{speckle} están
correlacionados dentro de cada grano --de tamaño transversal $0.378\lambda$ en
$\tilde z=1\lambda$, es decir unas $19$ muestras por grano en esta malla--, lo
que reduce drásticamente el número efectivo de muestras independientes y sesga
el $p$-valor hacia el rechazo. Por ello el test se reporta aquí como estadístico
descriptivo y no como criterio de aceptación. La fracción de píxeles con
$I>2\langle I\rangle$, que no depende de esa hipótesis, se sitúa entre $0.1299$
y $0.1543$ frente al valor teórico $e^{-2}=0.1353$, consistente con la
distribución esperada.

% ──────────────────────────────────────────────────────────────────────────────
\section{Contexto en la literatura}

La Tabla~\ref{tab:comparativa} sitúa los resultados de precisión de la
PINN-SIREN propuesta junto al antecedente más cercano en la literatura.

\begin{table}[htbp]
\centering
\begin{threeparttable}
\caption{Error $L^2$ relativo: PINN-SIREN y antecedente de la literatura.}
\label{tab:comparativa}
\begin{tabular}{lccc}
\toprule
Método & Dimensión & Comparación & Error $L^2$ (\%) \\
\midrule
\textcite{schoder2024helmholtz} & 3D (acústica) & PINN vs.\ FEM (malla fina) & 2.49 \\
\textcite{zhang2025fepirbn} & 2D (EM) & PINN vs.\ simulación numérica & 1.40--5.82 \\
\textbf{PINN-SIREN (este trabajo)} & \textbf{1D (óptica)} & PINN vs.\ solución analítica & \textbf{0.006} \\
\textbf{PINN-SIREN (este trabajo)} & \textbf{2D (óptica)} & PINN vs.\ solución analítica & \textbf{0.171} \\
\textbf{PINN-SIREN modal (este trabajo)} & \textbf{2D (óptica)} & PINN vs.\ espectro angular & \textbf{3.161}\tnote{a} \\
\bottomrule
\end{tabular}
\begin{tablenotes}\small
    \item Los resultados no son directamente comparables entre sí: difieren en
    dimensionalidad, dominio físico (acústico/electromagnético vs.\ óptico) y
    tipo de solución de referencia (FEM/simulación numérica vs.\ analítica o
    semianalítica). Se presentan como contexto de la literatura, no como línea
    base de comparación cuantitativa. Los resultados propios se validan
    contra el umbral fijo de la tesis ($L^2 < 5\%$), independiente de
    trabajos externos.
    \item[a] Valor de campo completo. Está dominado por la fracción evanescente
    de la referencia, irrepresentable por la base modal: el error sobre el
    subespacio propagante es de $0.042\%$ en promedio
    (Sección~\ref{sec:nb03_resultados_z1}). La comparación con los dos trabajos
    externos de esta tabla es además de distinta naturaleza, porque ambos
    contrastan contra un método numérico independiente y aquí la referencia es
    semianalítica (Sección~\ref{sec:nb03_referencia}).
\end{tablenotes}
\end{threeparttable}
\end{table}

El antecedente más cercano en arquitectura a este trabajo es
\textcite{panagiotakopoulos2026helmholtz}, quienes resuelven la misma ecuación
de Helmholtz 2D con campo complejo, usando también una red SIREN y la misma
estrategia de optimización bifásica Adam seguido de L-BFGS. Se diferencia de
este trabajo en tres aspectos: (i) validan cualitativamente la física del
campo resultante (propagación, refracción, dispersión) en lugar de reportar
un error $L^2$ cuantitativo contra una solución analítica; (ii) no calibran
$\omega_0$ respecto al número de onda $k$; y (iii) su condición de frontera es
una fuente puntual gaussiana con condición de radiación de Sommerfeld, en
lugar de la fase aleatoria en $z=0$ que da lugar al patrón de \textit{speckle}
en este trabajo. No aborda, por tanto, la generación de \textit{speckle} ni su
validación estadística contra los criterios de \textcite{goodman2007speckle}.

% ──────────────────────────────────────────────────────────────────────────────
\section{Resumen de experimentos completados}

La Tabla~\ref{tab:resumen_experimentos} sintetiza el estado de los experimentos
planificados en los objetivos específicos.

\begin{table}[H]
\centering
\begin{threeparttable}
\caption{Estado de los experimentos planificados.}
\label{tab:resumen_experimentos}
\begin{tabular}{llcc}
\toprule
Experimento & Descripción & Estado & Métrica principal \\
\midrule
NB01 & Helmholtz 1D, validación analítica & Completo & $L^2 = 0.006\%$ \\
NB02 & Helmholtz 2D, campo complejo       & Completo & $L^2 = 0.171\%$ \\
NB03-A & Speckle, referencia por espectro angular & Completo & $C=0.9661$ (conjunto, 64 real.) \\
NB03-B & PINN-SIREN modal en $\tilde z=1\lambda$ & Completo & $L^2 = 3.161\%$ (5 pantallas) \\
NB04 & Benchmark FEM, factor de aceleración & Pendiente & $S = T_\text{FEM}/T_\text{PINN}$ \\
\bottomrule
\end{tabular}
\begin{tablenotes}\small
    \item Los estados se refieren al cierre de esta versión de la tesis. El
    contraste de NB03-A es el promedio del contraste de conjunto de las cinco
    pantallas, cada uno sobre $64$ realizaciones en $\tilde z=1\lambda$. NB04
    permanece abierto para trabajo posterior y es el único objetivo específico
    sin evidencia asociada.
\end{tablenotes}
\end{threeparttable}
\end{table}

Los tres experimentos completados verifican la hipótesis de precisión. NB01 y
NB02 cumplen el umbral $L^2 < 5\%$ con amplitud frente a soluciones analíticas.
NB03 lo cumple en las cinco pantallas de fase independientes frente a la
referencia por espectro angular, con la precisión adicional de que el error
residual está dominado por la fracción evanescente del espectro y no por el
ajuste de la red (Sección~\ref{sec:nb03_resultados_z1}). La hipótesis
estadística se cumple en los dos niveles definidos en
Sección~\ref{sec:nb03_goodman}. La hipótesis de velocidad ($S > 1$) permanece
sin responder: requiere la implementación de la referencia FEM (NB04), que es
además la que aportaría una referencia numérica independiente del método, según
se discute en la Sección~\ref{sec:nb03_referencia}.
